% =====================================================================
%  Gemeinsamer Inhalt der Betriebsanweisungen für die Bambu-Lab-Drucker
%  (wird von BA_X1C.tex, BA_P1S.tex und Einweisung_3D-Drucker.tex eingebunden)
% =====================================================================

\BAkopf{Drucken, Filamentwechsel, Teileentnahme, Reinigung}

% ---------- 1 Anwendungsbereich --------------------------------------
\BAblock{ANWENDUNGSBEREICH}{M002}{%
  \begin{itemize}
    \item Additive Fertigung von Kunststoffteilen im FDM-Verfahren.
    \item Zugelassene Materialien: \textbf{nur PLA und ASA}. Andere
          Filamente nur nach Freigabe durch die FabLab-Betreuer.
    \item Benutzung nur nach Einweisung; Hinweise des Herstellers beachten.
  \end{itemize}}

% ---------- 2 Gefahren -----------------------------------------------
\BAblock{GEFAHREN FÜR MENSCH UND UMWELT}{W017,W024,W012}{%
  \begin{itemize}
    \item \textbf{Verbrennungsgefahr} an Düse/Hotend (bis 300\,°C)
          und Druckbett (bis ca.\ 100\,°C bei ASA).
    \item \textbf{Quetschgefahr} durch schnell bewegte Achsen bei geöffneter Tür.
    \item \textbf{Emissionen} beim Druck von \textbf{ASA}: Styrol und weitere
          flüchtige organische Verbindungen (VOC) sowie ultrafeine Partikel;
          Geruchsbelästigung möglich. PLA ist emissionsarm.
    \item \textbf{Schnittgefahr} an scharfkantigen Druckteilen und Stützstrukturen.
    \item \textbf{Elektrische Gefährdung} bei beschädigtem Netzkabel oder
          geöffnetem Gehäuse (230\,V).
    \item \textbf{Brandgefahr} bei Fehlfunktion (Filamentstau,
          Überhitzung, Verunreinigungen im Bauraum).
  \end{itemize}}

% ---------- 3 Schutzmaßnahmen ----------------------------------------
\BAblock{SCHUTZMASSNAHMEN UND VERHALTENSREGELN}{P015}{%
  \begin{itemize}
    \item Tür und Deckel während des Drucks geschlossen halten.
    \item Nicht in den Bauraum greifen, solange sich Achsen bewegen.
    \item Nach dem Druck Druckbett abkühlen lassen; Druckplatte nur am
          Rand anfassen und herausnehmen.
    \item Teile nur durch Biegen der Druckplatte lösen, keine Spachtel oder
          Messer (PEI nicht verkratzen).
    \item Nach ASA-Drucken Bauraum einige Minuten geschlossen lassen
          (Filter nimmt Dämpfe auf).
    \item Lüftungsöffnungen freihalten; keine brennbaren Gegenstände
          auf oder neben dem Drucker ablegen.
    \item Lange Druckjobs regelmäßig kontrollieren.
    \item Essen und Trinken am Drucker verboten.
  \end{itemize}}

% ---------- 4 Störungen ----------------------------------------------
\BAblock{VERHALTEN BEI STÖRUNGEN}{W001,F001}{%
  \begin{itemize}
    \item Druck über das Display oder Bambu Studio stoppen.
    \item Bei Rauch, Brandgeruch oder Funkenbildung sofort Drucker
          ausschalten (Schalter hinten rechts) oder \textbf{Netzstecker ziehen}.
    \item Störungen (z.\,B. Filamentstau, verstopfte Düse) nicht selbst
          beheben, sondern den FabLab-Betreuern melden.
    \item Beschädigte Geräte nicht weiter betreiben und kennzeichnen.
    \item Entstehungsbrand: Gerät stromlos machen, mit CO\textsubscript{2}-Löscher
          löschen. \textbf{Feuerwehr: 112}
  \end{itemize}}

% ---------- 5 Erste Hilfe --------------------------------------------
\BAblock{VERHALTEN BEI UNFÄLLEN – ERSTE HILFE}{E003}{%
  \begin{itemize}
    \item Ruhe bewahren, Gerät abschalten, Betreuer/Ersthelfer verständigen.
    \item Verbrennungen: sofort 10--15 Min.\ mit handwarmem Wasser kühlen,
          keine Salben; bei großflächigen Verbrennungen Arzt aufsuchen.
    \item Bei Atembeschwerden: an die frische Luft, ärztliche Hilfe holen.
    \item Jeden Unfall den Betreuern melden und im Verbandbuch eintragen.
  \end{itemize}\vspace{1.5mm}
  \textbf{Notruf: 112}\qquad FabLab: 09131\,/\,85\,28013}

% ---------- 6 Instandhaltung / Entsorgung (mit Fußzeile) --------------
\BAletzterblock{INSTANDHALTUNG, ENTSORGUNG}{M006}{%
  \begin{itemize}
    \item Wartung (Düse, Achsen, Aktivkohlefilter) nur durch
          eingewiesene Betreuer im ausgeschalteten, abgekühlten Zustand.
    \item Elektrische Prüfung nach DGUV Vorschrift 3 gemäß festgelegter
          Prüffrist.
    \item Filamentreste, Fehldrucke und Stützmaterial im Restmüll entsorgen.
  \end{itemize}}
